% \begin{theorem}[Fubini Theorem for Colimits and Limtis] \label{theorem:colimits_commute_with_colimits_and_limits_commute_with_limits}
%     Let $\mathcal{C}$ be a category and let $I$ and $J$ be small categories. Let $D: I \times J \to \mathcal{C}$ be a bifunctor. 
    
%     \begin{enumerate}
%         \item Assume that the following \CrefAndHyperrefIfExist{definition:limit_and_colimit_of_a_diagram_in_a_category}{colimits} exist in $\calC$:
%         \begin{itemize}
%             \item for every object $i \in I$, $\varinjlim_{j \in J} D(i, j)$
%             \item $\varinjlim_{i \in I} \left( \varinjlim_{j \in J} D(i, j) \right)$
%             \item for every object $j \in J$, $\varinjlim_{i \in I} D(i, j)$
%             \item $\varinjlim_{j \in J} \left( \varinjlim_{i \in I} D(i, j) \right)$.
%         \end{itemize}
%         Then $\varinjlim_{(i, j) \in I \times J} D(i, j)$ exists and there are canonical isomorphisms:
%         $$
%         \varinjlim_{i \in I} \left( \varinjlim_{j \in J} D(i, j) \right) 
%         \cong \varinjlim_{(i, j) \in I \times J} D(i, j) 
%         \cong \varinjlim_{j \in J} \left( \varinjlim_{i \in I} D(i, j) \right).
%         $$

%                 \item Assume that the following \CrefAndHyperrefIfExist{definition:limit_and_colimit_of_a_diagram_in_a_category}{limits} exist in $\calC$:
%         \begin{itemize}
%             \item for every object $i \in I$, $\varprojlim_{j \in J} D(i, j)$
%             \item $\varprojlim_{i \in I} \left( \varprojlim_{j \in J} D(i, j) \right)$
%             \item for every object $j \in J$, $\varprojlim_{i \in I} D(i, j)$
%             \item $\varprojlim_{j \in J} \left( \varprojlim_{i \in I} D(i, j) \right)$.
%         \end{itemize}
%         Then $\varprojlim_{(i, j) \in I \times J} D(i, j)$ exists and there are canonical isomorphisms:
%         $$
%         \varprojlim_{i \in I} \left( \varprojlim_{j \in J} D(i, j) \right) 
%         \cong \varprojlim_{(i, j) \in I \times J} D(i, j) 
%         \cong \varprojlim_{j \in J} \left( \varprojlim_{i \in I} D(i, j) \right).
%         $$

%     \end{enumerate}
% \end{theorem}
\begin{theorem}[Fubini Theorem for Colimits and Limits] \label{theorem:colimits_commute_with_colimits_and_limits_commute_with_limits}
    Let $\mathcal{C}$ be a category and let $I$ and $J$ be small categories. Let $D: I \times J \to \mathcal{C}$ be a bifunctor. 
    
    \begin{enumerate}
        \item Assume the following \CrefAndHyperrefIfExist{definition:limit_and_colimit_of_a_diagram_in_a_category}{colimits} exist in $\calC$:
        \begin{itemize}
            \item for every object $i \in I$, $\varinjlim_{j \in J} D(i, j)$
            \item $\varinjlim_{(i, j) \in I \times J} D(i, j)$
            \item $\varinjlim_{i \in I} (\varinjlim_{j \in J} D(i, j))$
        \end{itemize}
        There is a canonical isomorphism:
        $$
        \varinjlim_{i \in I} \left( \varinjlim_{j \in J} D(i, j) \right) \cong \varinjlim_{(i, j) \in I \times J} D(i, j).
        $$

        \item Assume the following \CrefAndHyperrefIfExist{definition:limit_and_colimit_of_a_diagram_in_a_category}{limits} exist in $\calC$:
        \begin{itemize}
            \item for every object $i \in I$, $\varprojlim_{j \in J} D(i, j)$
            \item $\varprojlim_{(i, j) \in I \times J} D(i, j)$
            \item $\varprojlim_{i \in I} (\varprojlim_{j \in J} D(i, j))$
        \end{itemize}
        There is a canonical isomorphism:
        $$
        \varprojlim_{i \in I} \left( \varprojlim_{j \in J} D(i, j) \right) \cong \varprojlim_{(i, j) \in I \times J} D(i, j).
        $$
    \end{enumerate}
\end{theorem}
\begin{remark}
    These results assert that the \CrefAndHyperrefIfExist{theorem:limits_and_colimits_as_functors_from_functor_category_to_value_category}{colimit functor} $\operatorname{colim}: \mathcal{C}^{I \times J} \to \mathcal{C}$ (assuming $\mathcal{C}$ is cocomplete) is isomorphic to the composition of partial colimit functors $\operatorname{colim}_I \circ \operatorname{colim}_J$ and $\operatorname{colim}_J \circ \operatorname{colim}_I$.
\end{remark}
\begin{proof}
    \TODO{AI generated, verify}
    We prove the colimit version. The limit version follows dually.

    Let $L = \varprojlim_{(i, j) \in I \times J} D(i, j)$ denote the limit over the product category. We want to show that $L' = \varprojlim_{i \in I} (\varprojlim_{j \in J} D(i, j))$ is isomorphic to $L$.

    The limit $L$ comes with canonical projections $\pi_{i,j} : L \to D(i, j)$ for every pair $(i, j)$.
    
    \begin{enumerate}
        \item \textbf{Constructing a map $\phi: L \to L'$:}
        For a fixed $i \in I$, the collection of projections $\{\pi_{i,j} : L \to D(i, j)\}_{j \in J}$ forms a cone over the diagram $D(i, -)$. By the universal property of $L_i = \varprojlim_{j \in J} D(i, j)$, there exists a unique morphism $\alpha_i : L \to L_i$ such that for all $j \in J$, the composition $L \xrightarrow{\alpha_i} L_i \to D(i, j)$ is $\pi_{i,j}$.
        
        The collection $\{\alpha_i : L \to L_i\}_{i \in I}$ now forms a cone over the diagram $i \mapsto L_i$. By the universal property of $L' = \varprojlim_{i \in I} L_i$, there exists a unique morphism $\phi: L \to L'$ such that for all $i \in I$, the composition $L \xrightarrow{\phi} L' \to L_i$ is $\alpha_i$.

        \item \textbf{Constructing a map $\psi: L' \to L$:}
        Let $p_i : L' \to L_i$ be the projection from the limit $L'$, and $q_{i,j} : L_i \to D(i, j)$ be the projection from the limit $L_i$. The composition $q_{i,j} \circ p_i : L' \to D(i, j)$ constitutes a cone over the diagram $D$ indexed by $I \times J$ (i.e., for any morphism $(u, v) : (i, j) \to (i', j')$, the corresponding diagrams commute). By the universal property of $L = \varprojlim_{I \times J} D$, there exists a unique morphism $\psi: L' \to L$ such that $\pi_{i,j} \circ \psi = q_{i,j} \circ p_i$.

        \item \textbf{Verification:}
        One verifies that $\psi \circ \phi = \text{id}_L$ and $\phi \circ \psi = \text{id}_{L'}$ by checking that these compositions induce the identity on the projections $\pi_{i,j}$ and $p_i$ respectively, invoking the uniqueness part of the universal properties.
    \end{enumerate}
\end{proof}
